\chapter{Prompt engineering}
\label{ch:prompt-engineering}

\begin{quote}
\emph{<< Le prompt engineering est l'art de communiquer avec un mod\`ele de langage. Les capacit\'es du mod\`ele sont fix\'ees ; votre prompt d\'etermine combien de ces capacit\'es vous d\'everrouillez. >>}
\end{quote}

% --- hook ---
Un soir de 2022, Jason Wei et ses collègues de Google découvrent qu'ajouter cinq mots à la fin d'un prompt --- \emph{Let's think step by step} --- fait passer la précision de GPT-3 sur les problèmes mathématiques de 18\% à 79\%. Aucun changement de modèle, aucun fine-tuning, juste ces cinq mots. C'est l'instant où le \emph{prompt engineering} cesse d'être de la bricole pour devenir une discipline. Les capacités du modèle sont fixes ; ce que vous en obtenez l'est beaucoup moins. Ce chapitre n'enseigne pas des astuces : il donne une boîte à outils structurée --- zero-shot, few-shot, chain-of-thought, sortie structurée, prompt par rôle --- avec les cas où chaque outil fonctionne et ceux où il ne fonctionne pas.
\section{Pourquoi le prompt compte}

Un grand mod\`ele de langage est un pr\'edicteur de texte g\'en\'eraliste. Le prompt cadre la t\^ache, fournit le contexte et contraint le format de sortie. Le m\^eme mod\`ele peut \^etre traducteur, g\'en\'erateur de code ou assistant m\'edical --- enti\`erement selon le prompt.

% ============================================================
\section{Mise en place : acc\`es API gratuit}

Nous utilisons l'API gratuite Groq (inf\'erence rapide sur Llama 3 et Mixtral) et le palier gratuit de Google Gemini :

\begin{minted}{python}
# Installation des clients
# !pip install groq google-generativeai

import os
# D\'efinir vos cl\'es API (cl\'es gratuites depuis console.groq.com et aistudio.google.com)
os.environ["GROQ_API_KEY"] = "your-groq-key"
os.environ["GOOGLE_API_KEY"] = "your-google-key"
\end{minted}

\begin{minted}{python}
from groq import Groq

client = Groq()

def ask_groq(prompt, model="llama-3.1-8b-instant", temperature=0.7):
    """Envoie un prompt \`a Groq et renvoie la r\'eponse."""
    response = client.chat.completions.create(
        model=model,
        messages=[{"role": "user", "content": prompt}],
        temperature=temperature,
        max_tokens=1024,
    )
    return response.choices[0].message.content

print(ask_groq("What is the capital of Benin?"))
\end{minted}

\begin{minted}{python}
import google.generativeai as genai

genai.configure()
gemini = genai.GenerativeModel("gemini-2.0-flash")

def ask_gemini(prompt, temperature=0.7):
    """Envoie un prompt \`a Gemini et renvoie la r\'eponse."""
    response = gemini.generate_content(
        prompt,
        generation_config=genai.GenerationConfig(temperature=temperature),
    )
    return response.text

print(ask_gemini("What is the capital of Benin?"))
\end{minted}

% ============================================================
\section{Prompt zero-shot}

Donner au mod\`ele une t\^ache sans aucun exemple :

\begin{minted}{python}
prompt = """Classify the following movie review as POSITIVE or NEGATIVE.

Review: "The cinematography was breathtaking, but the plot was
completely incoherent and the acting felt wooden."

Classification:"""

print(ask_groq(prompt, temperature=0))
\end{minted}

% ============================================================
\section{Prompt few-shot}

Fournir des exemples pour guider le comportement du mod\`ele :

\begin{minted}{python}
prompt = """Translate English to French.

English: The weather is beautiful today.
French: Le temps est magnifique aujourd'hui.

English: I would like a cup of coffee, please.
French: Je voudrais une tasse de cafe, s'il vous plait.

English: Where is the nearest hospital?
French:"""

print(ask_groq(prompt, temperature=0.2))
\end{minted}

\begin{aitip}
Les exemples few-shot doivent \^etre vari\'es, de haute qualit\'e, et repr\'esentatifs des entr\'ees attendues. 3 \`a 5 exemples suffisent g\'en\'eralement. Plus d'exemples peut am\'eliorer la pr\'ecision sur des t\^aches structur\'ees.
\end{aitip}

% ============================================================
\section{Chain-of-thought (CoT)}

Demander au mod\`ele de raisonner \'etape par \'etape avant de donner une r\'eponse finale :

\begin{minted}{python}
# Sans CoT
prompt_direct = "If a train travels 120 km in 1.5 hours, and then 80 km in 1 hour, what is the average speed for the entire trip?"

# Avec CoT
prompt_cot = """If a train travels 120 km in 1.5 hours, and then 80 km in 1 hour, what is the average speed for the entire trip?

Let's think step by step:"""

print("=== Direct ===")
print(ask_groq(prompt_direct, temperature=0))
print("\n=== Chain-of-Thought ===")
print(ask_groq(prompt_cot, temperature=0))
\end{minted}

Le chain-of-thought am\'eliore consid\'erablement la pr\'ecision sur les t\^aches math\'ematiques, de logique et de raisonnement multi-\'etapes.

% ============================================================
\section{Prompt par r\^ole}

Assigner un r\^ole au mod\`ele via le message syst\`eme :

\begin{minted}{python}
from groq import Groq
client = Groq()

response = client.chat.completions.create(
    model="llama-3.1-8b-instant",
    messages=[
        {"role": "system", "content": (
            "You are an expert data scientist. You explain concepts "
            "clearly using analogies. You always provide Python code "
            "examples. You never use jargon without defining it first."
        )},
        {"role": "user", "content": "Explain overfitting to a beginner."},
    ],
    temperature=0.7,
)
print(response.choices[0].message.content)
\end{minted}

% ============================================================
\section{Sortie structur\'ee}

Forcer le mod\`ele \`a produire du JSON, CSV ou d'autres formats structur\'es :

\begin{minted}{python}
prompt = """Extract the following information from the text and return
it as a JSON object with keys: name, age, condition, medication.

Text: "Mrs. Fatou Diallo, 67 years old, was diagnosed with
Type 2 diabetes last year. She takes Metformin 500mg twice daily."

JSON:"""

import json
result = ask_groq(prompt, temperature=0)
print(result)
data = json.loads(result)
print(f"Patient: {data['name']}, Age: {data['age']}")
\end{minted}

\begin{minted}{python}
# Avec Gemini et la configuration de sortie structur\'ee
import google.generativeai as genai
import typing_extensions as typing

class PatientInfo(typing.TypedDict):
    name: str
    age: int
    condition: str
    medication: str

model = genai.GenerativeModel("gemini-2.0-flash")
result = model.generate_content(
    "Extract patient info: Mrs. Fatou Diallo, 67, Type 2 diabetes, Metformin 500mg",
    generation_config=genai.GenerationConfig(
        response_mime_type="application/json",
        response_schema=PatientInfo,
    ),
)
print(result.text)
\end{minted}

% ============================================================
\section{Templates de prompts}

Construire des prompts r\'eutilisables avec variables :

\begin{minted}{python}
from string import Template

summarizer = Template("""Summarize the following $doc_type in $num_sentences sentences.
Focus on: $focus_area

Text:
$text

Summary:""")

prompt = summarizer.substitute(
    doc_type="research paper abstract",
    num_sentences="3",
    focus_area="methodology and key findings",
    text="We propose LoRA, a method for adapting large language models..."
)
print(ask_groq(prompt))
\end{minted}

\begin{minted}{python}
# Avec les templates de prompts LangChain
from langchain_core.prompts import ChatPromptTemplate

template = ChatPromptTemplate.from_messages([
    ("system", "You are a {role}. Respond in {language}."),
    ("human", "{question}"),
])

prompt = template.invoke({
    "role": "medical doctor",
    "language": "French",
    "question": "What are the symptoms of malaria?",
})
print(prompt.to_string())
\end{minted}

% ============================================================
\section{Patterns de prompts courants}

\begin{longtable}{p{3.5cm}p{9cm}}
\toprule
\textbf{Pattern} & \textbf{Exemple} \\
\midrule
Persona & << You are a senior Python developer... >> \\
Cadrage de t\^ache & << Your task is to classify emails as spam or not-spam. >> \\
Format de sortie & << Return your answer as a JSON array. >> \\
Contraintes & << Use only information from the provided text. >> \\
\'Etape par \'etape & << Think step by step before answering. >> \\
Self-consistency & Ex\'ecuter 3 chemins CoT, prendre la r\'eponse majoritaire. \\
\bottomrule
\end{longtable}

\begin{exercise}
\begin{enumerate}
  \item \'Ecrivez un prompt zero-shot qui classe des tickets de support client en cat\'egories : facturation, technique, compte, g\'en\'eral. Testez sur 5 exemples.
  \item Cr\'eez un prompt few-shot (3 exemples) pour la reconnaissance d'entit\'es nomm\'ees. Le mod\`ele doit extraire noms de personnes, organisations et lieux d'un texte.
  \item Utilisez le chain-of-thought pour r\'esoudre : << Un magasin vend des pommes \`a 1,50\$ chacune. Un client en ach\`ete 3 et paie avec un billet de 10\$. Il a aussi un coupon de 20\% de remise. Quelle est sa monnaie ? >>
  \item Construisez un template de prompt r\'eutilisable pour la revue de code. Le template prend en entr\'ee : langage, extrait de code, et focus de revue (s\'ecurit\'e, performance, lisibilit\'e).
  \item Comparez le m\^eme prompt sur Groq (Llama 3) et Gemini. Notez les diff\'erences de format, verbosit\'e et pr\'ecision.
\end{enumerate}
\end{exercise}

\begin{ethicsbox}
Le prompt engineering peut servir \`a contourner les garde-fous de s\^uret\'e (<< jailbreaking >>). En tant que praticiens de l'IA, nous avons la responsabilit\'e de ne pas d\'evelopper ni partager de prompts de jailbreak. Pour tester la s\^uret\'e d'un mod\`ele, utilisez des environnements contr\^ol\'es et signalez les vuln\'erabilit\'es aux fournisseurs via une divulgation responsable.
\end{ethicsbox}

% ============================================================
\section{R\'esum\'e du chapitre}

\begin{itemize}
  \item Le prompt est l'interface entre l'intention humaine et le comportement du mod\`ele.
  \item Le zero-shot fonctionne pour les t\^aches simples ; le few-shot ajoute des exemples pour les t\^aches plus difficiles.
  \item Le chain-of-thought am\'eliore consid\'erablement le raisonnement en demandant au mod\`ele de montrer son travail.
  \item Le prompt par r\^ole via les messages syst\`eme contr\^ole ton, expertise et format.
  \item La sortie structur\'ee (JSON, sch\'emas) rend les r\'eponses LLM exploitables par machine.
  \item Les templates de prompts rendent les prompts r\'eutilisables et param\'etrables.
  \item Les API gratuites (Groq, Gemini) donnent acc\`es \`a des mod\`eles de pointe sans co\^ut.
\end{itemize}
